\lampiransection{Hasil Wawancara Awal Pemilik Nasi Gerilya}
\label{app:wawancara-awal-pemilik}

Lampiran ini memuat hasil wawancara awal dengan pemilik usaha untuk penelitian \textit{Strategi Pemasaran Rumah Makan Nasi Gerilya pada Platform GrabFood untuk Meningkatkan Penjualan Berdasarkan Analisis SWOT}. Wawancara awal dilakukan untuk memperoleh data mengenai profil usaha, penggunaan platform, struktur pelanggan, kekuatan usaha, strategi pemasaran, kendala operasional, peluang pengembangan, dan ancaman persaingan. Seluruh poin pada lampiran ini disusun berdasarkan data dan keterangan yang disampaikan oleh pemilik usaha pada saat wawancara awal.

\begingroup
\begin{xltabular}{\textwidth}{|Z{3.4cm}|Y|}
\caption{Identitas Wawancara Awal}
\label{tab:identitas_wawancara_awal} \\
        \toprule
        \textbf{Aspek}             & \textbf{Keterangan}                                                                                       \\ \hline
        \midrule
\endfirsthead
\continuedtabletitle{2}
        \toprule
        \textbf{Aspek}             & \textbf{Keterangan}                                                                                       \\ \hline
        \midrule
\endhead
\continuedtablefooter{2}
\endfoot
        Jenis kegiatan             & Wawancara awal pra-survei                                                                                 \\ \hline
        Waktu pelaksanaan          & April 2026                                                                                                \\ \hline
        Informan                   & Pemilik atau pengelola Rumah Makan Nasi Gerilya                                                           \\ \hline
        Fokus wawancara            & Profil usaha, kanal penjualan, pelanggan, strategi pemasaran, kendala operasional, peluang, dan ancaman   \\ \hline
        Bentuk data                & Catatan wawancara yang diringkas dan diolah peneliti tanpa mengubah substansi keterangan informan         \\ \hline
        Fungsi dalam penelitian    & Data awal untuk menyusun pedoman wawancara utama dan memperkuat identifikasi faktor internal-eksternal    \\ \hline
        \bottomrule
    
\end{xltabular}
\tablesource[\textwidth]{Wawancara awal dengan pemilik usaha; diolah peneliti (2026).}
\endgroup

Tabel berikut menyajikan daftar pertanyaan pra-survei dan jawaban pemilik usaha berdasarkan hasil wawancara awal.

    {\scriptsize
        

\refstepcounter{table}\label{tab:wawancara_awal_pemilik}
        \edef\WawancaraAwalNomor{\thetable}
        \addcontentsline{lot}{table}{\protect\numberline{\thetable}Pertanyaan dan Jawaban Pra-Survei Pemilik Rumah Makan Nasi Gerilya}
        \begin{xltabular}{\textwidth}{|>{\centering\arraybackslash}p{0.7cm}|p{4.4cm}|Y|}
            \multicolumn{3}{c}{\textbf{Tabel \WawancaraAwalNomor}}                                                                                                                                                                                                                                                                                                                                                                                                                                  \\
            \multicolumn{3}{c}{\textbf{Pertanyaan dan Jawaban Pra-Survei Pemilik Rumah Makan Nasi Gerilya}}                                                                                                                                                                                                                                                                                                                                                                                         \\[6pt]
            \toprule
            \textbf{No.} & \textbf{Pertanyaan}                                                                                                            & \textbf{Jawaban Pemilik}                                                                                                                                                                                                                                                                                                                      \\ \hline
            \midrule
            \endfirsthead
            \continuedtabletitle{3}
            \toprule
            \textbf{No.} & \textbf{Pertanyaan}                                                                                                            & \textbf{Jawaban Pemilik}                                                                                                                                                                                                                                                                                                                      \\ \hline
            \midrule
            \endhead
            \continuedtablefooter{3}
            \endfoot
            \bottomrule
            \endlastfoot
            1            & Sejak kapan Rumah Makan Nasi Gerilya mulai berdiri dan mulai bergabung dengan platform GrabFood?                               & Nasi Gerilya mulai jalan sejak September 2024. Dari awal kami memang sudah masuk ke kanal online, termasuk GrabFood, karena targetnya bukan hanya pelanggan yang datang langsung, tetapi juga pelanggan yang pesan dari aplikasi.                                                                                                             \\ \hline
            2            & Kanal penjualan online apa saja yang digunakan Rumah Makan Nasi Gerilya saat ini?                                              & Sekarang kami pakai GrabFood, GoFood, dan ShopeeFood. Tapi dari tiga platform itu, GrabFood yang paling kami fokuskan karena paling cocok dengan pasar kami.                                                                                                                                                                                  \\ \hline
            3            & Apa alasan utama usaha ini memasarkan produknya melalui GrabFood?                                                              & GrabFood itu paling cocok dengan pasar kami. Dari beberapa platform, GrabFood terasa paling sesuai dengan segmen pelanggan Nasi Gerilya dan kontribusinya juga paling besar ke penjualan, diperkirakan lebih dari 50\% pendapatan penjualan. Jadi bukan sekadar tambahan, tapi memang jadi kanal utama.                                       \\ \hline
            4            & Bagaimana pandangan Bapak/Ibu terhadap GoFood dan ShopeeFood dibandingkan GrabFood?                                            & GoFood menurut kami penggunanya lebih banyak tersebar di pinggiran kota, jadi kurang pas dengan segmen yang kami incar. ShopeeFood lebih kuat di pembeli yang mencari voucher, tapi belum tentu mereka jadi pelanggan ulang.                                                                                                                  \\ \hline
            5            & Bagaimana komposisi pelanggan online Rumah Makan Nasi Gerilya saat ini?                                                        & Untuk pelanggan online, saat ini kira-kira 55\% pelanggan baru, 38\% repeat order, dan 7\% reactivated user. Sebelumnya pelanggan baru sekitar 45\%, repeat order 43\%, dan reactivated user 12\%. Perubahan itu kami kaitkan juga dengan marketing plan high season.                                                                         \\ \hline
            6            & Menurut Bapak/Ibu, apa kekuatan utama Rumah Makan Nasi Gerilya dibanding rumah makan lain di GrabFood?                         & Kekuatan kami ada di rasa yang otentik dan karakter brand yang cukup kuat. Kami ingin orang ingat Nasi Gerilya sebagai nasi Padang yang rasanya khas, bukan hanya ikut jualan di aplikasi.                                                                                                                                                    \\ \hline
            7            & Menu apa yang paling diunggulkan atau paling sering dipesan pelanggan melalui GrabFood?                                        & Yang paling kami unggulkan itu dendeng dan ayam pop putih. Dari tampilan GrabFood juga menu seperti rendang sapi, ayam goreng, ayam bakar, dan kikil tunjang termasuk menu yang sering dilihat dan dipesan pelanggan.                                                                                                                         \\ \hline
            8            & Bagaimana strategi promosi yang selama ini digunakan untuk menarik pelanggan di GrabFood?                                      & Promosi kami tidak hanya mengandalkan GrabFood. GrabFood lebih sering jadi tempat pelanggan langsung order. Untuk menarik perhatian, kami banyak dorong dari media sosial, KOL, influencer, promo platform, dan penguatan tampilan produk.                                                                                                    \\ \hline
            9            & Apakah promo seperti diskon, voucher, atau potongan harga cukup berpengaruh terhadap penjualan?                                & Berpengaruh, apalagi kalau ada voucher atau program GrabFood. Tapi sekarang promo juga harus dihitung baik-baik karena tidak semuanya ditanggung platform. Ada biaya, pajak, dan margin yang tetap harus dijaga.                                                                                                                              \\ \hline
            10           & Apakah ada pertimbangan khusus terkait harga, margin, dan biaya promosi di platform online?                                    & Ada. Harga memang harus dihitung ulang karena kami juga sedang menjaga margin dan net profit. Selain itu, biaya promo, pajak, dan diskon sekarang tidak selalu ditanggung platform, jadi harus benar-benar dihitung supaya penjualan tetap sehat.                                                                                             \\ \hline
            11           & Menurut Bapak/Ibu, apa kelemahan atau kendala utama yang dihadapi dalam pemasaran melalui GrabFood?                            & Kendalanya cukup banyak. Ongkir kadang terasa mahal, jarak layanan terbatas, dan pelanggan sensitif dengan harga. Dari operasional, masalahnya ada di menu yang habis saat jam ramai, antrean driver, dan konsistensi rasa yang masih perlu dijaga.                                                                                           \\ \hline
            12           & Kendala operasional apa yang paling sering muncul saat pesanan online sedang ramai?                                            & Saat peak hour, beberapa menu bisa sold out dan perlu restock. Kadang menu yang habis masih bisa dipesan, lalu pelanggan bisa batal. Kapasitas penggorengan juga masih kurang, antrean driver bisa menumpuk, dan kami juga belum punya dessert untuk bantu upselling.                                                                         \\ \hline
            13           & Bagaimana kondisi konsistensi rasa dan proses kerja dapur saat ini?                                                            & Konsistensi rasa masih jadi perhatian besar. Ada cook yang kadang ambil shortcut, jadi perlu pengawasan lebih ketat supaya rasa tetap stabil dan pelanggan dapat pengalaman yang sama setiap pesan.                                                                                                                                           \\ \hline
            14           & Keluhan pelanggan seperti apa yang paling sering muncul terkait pesanan online?                                                & Biasanya soal pesanan yang tidak sesuai request, menu yang ternyata habis, atau rasa yang belum konsisten. Kadang saat peak hour dapur juga penuh, jadi ada risiko layanan tidak secepat atau serapi yang kami mau.                                                                                                                           \\ \hline
            15           & Bagaimana cara usaha menanggapi review atau ulasan dari pelanggan?                                                             & Review tetap kami perhatikan karena itu kelihatan langsung oleh pelanggan lain. Kalau ada keluhan, kami jadikan bahan evaluasi, terutama soal rasa, pesanan yang kurang sesuai, dan ketersediaan menu.                                                                                                                                        \\ \hline
            16           & Menurut Bapak/Ibu, peluang apa yang paling besar untuk meningkatkan penjualan melalui GrabFood?                                & Peluangnya ada di upselling dan product lineup. Misalnya menu bisa dibuat lebih lengkap, ada tambahan produk yang mendorong nilai transaksi, dan ikut program GrabFood seperti voucher, verified store, atau awards.                                                                                                                          \\ \hline
            17           & Ancaman apa yang paling dirasakan dari pesaing atau dari sistem platform GrabFood?                                             & Pesaing baru mulai banyak, apalagi brand yang sudah kuat seperti Pagi Sore atau Payakumbuah. Dari sisi platform, ancamannya ongkir, beban promo yang makin terasa ke pelanggan, dan pelanggan yang makin sensitif terhadap harga.                                                                                                             \\ \hline
            18           & Ke depan, strategi dan rencana pengembangan apa yang ingin dilakukan agar Rumah Makan Nasi Gerilya semakin unggul di GrabFood? & Kami ingin perkuat kualitas produk, layanan, dan konsistensi rasa dulu. Setelah itu, kami mau optimalkan product lineup, discount partnership, identitas brand seperti maskot atau signature gesture, lalu pelan-pelan buka cabang baru di Medan. Target jangka panjangnya bisa ekspansi ke Jakarta setelah cabang di Medan sudah lebih kuat. \\ \hline
        \end{xltabular}
        \tablesource[\textwidth]{Wawancara awal dengan pemilik usaha; diolah peneliti (2026).}
        \addtocounter{table}{-1}
    }


Data profil usaha, penggunaan platform, dan komposisi pelanggan berdasarkan wawancara awal disajikan dalam beberapa tabel berikut.

\begingroup
\begin{xltabular}{\textwidth}{|Z{3.4cm}|Y|}
\caption{Profil Umum Usaha Berdasarkan Wawancara Pemilik}
\label{tab:profil_wawancara_awal} \\
        \toprule
        \textbf{Aspek}                 & \textbf{Keterangan}              \\ \hline
        \midrule
\endfirsthead
\continuedtabletitle{2}
        \toprule
        \textbf{Aspek}                 & \textbf{Keterangan}              \\ \hline
        \midrule
\endhead
\continuedtablefooter{2}
\endfoot
        Nama usaha                     & Rumah Makan Nasi Gerilya         \\ \hline
        Tahun berdiri                  & September 2024                   \\ \hline
        Lama usaha saat wawancara awal & Sekitar 1 tahun 7 bulan          \\ \hline
        Kanal penjualan online         & GrabFood, GoFood, dan ShopeeFood \\ \hline
        \bottomrule
    
\end{xltabular}
\tablesource[\textwidth]{Wawancara awal dengan pemilik usaha; diolah peneliti (2026).}
\endgroup

\begingroup
\begin{xltabular}{\textwidth}{|Z{3.4cm}|Y|}
\caption{Keterangan Pemilik Mengenai Penggunaan Platform}
\label{tab:wawancara_platform} \\
        \toprule
        \textbf{Aspek}                & \textbf{Keterangan Berdasarkan Wawancara}                                                                       \\ \hline
        \midrule
\endfirsthead
\continuedtabletitle{2}
        \toprule
        \textbf{Aspek}                & \textbf{Keterangan Berdasarkan Wawancara}                                                                       \\ \hline
        \midrule
\endhead
\continuedtablefooter{2}
\endfoot
        Platform utama                & GrabFood                                                                                                        \\ \hline
        Kontribusi terhadap penjualan & Disebut sebagai penyumbang penjualan terbesar dan diperkirakan menyumbang lebih dari 50\% pendapatan penjualan  \\ \hline
        Pertimbangan pemilik          & Segmen pengguna GrabFood dinilai lebih sesuai dengan target pasar usaha                                         \\ \hline
        Dukungan platform             & Awards, \textit{verified store}, dan \textit{voucher discount}                                                  \\ \hline
        Pandangan terhadap GoFood     & Dinilai lebih banyak memiliki sebaran pengguna di pinggiran kota dan kurang sesuai dengan segmen pasar usaha    \\ \hline
        Pandangan terhadap ShopeeFood & Dinilai lebih banyak mendorong pembelian berbasis voucher; pembelian tidak selalu berujung pada pembelian ulang \\ \hline
        \bottomrule
    
\end{xltabular}
\tablesource[\textwidth]{Wawancara awal dengan pemilik usaha; diolah peneliti (2026).}
\endgroup

\begingroup
\begin{xltabular}{\textwidth}{|Z{2.8cm}|Y|Y|}
\caption{Indikasi STP Awal Rumah Makan Nasi Gerilya}
\label{tab:stp_wawancara_awal} \\
        \toprule
        \textbf{Komponen}    & \textbf{Indikasi Awal}                                                                                                                                                                      & \textbf{Dasar Data}                                                                                                          \\ \hline
        \midrule
\endfirsthead
\continuedtabletitle{3}
        \toprule
        \textbf{Komponen}    & \textbf{Indikasi Awal}                                                                                                                                                                      & \textbf{Dasar Data}                                                                                                          \\ \hline
        \midrule
\endhead
\continuedtablefooter{3}
\endfoot
        Segmentasi           & Konsumen \textit{online food delivery} yang memesan makanan siap santap pada area layanan GrabFood, dengan komposisi pelanggan baru, pelanggan ulang, dan pelanggan yang diaktivasi kembali & Komposisi pelanggan menunjukkan 55\% \textit{new-comers}, 38\% \textit{repeat order}, dan 7\% \textit{reactivated user}      \\ \hline
        \textit{Targeting}   & GrabFood dipilih sebagai kanal utama karena dinilai paling sesuai dengan target pasar usaha dibanding GoFood dan ShopeeFood                                                                 & Pemilik menyebut GrabFood sebagai kanal penyumbang lebih dari 50\% pendapatan penjualan dan paling sesuai dengan segmen pasar usaha \\ \hline
        \textit{Positioning} & Nasi Gerilya diarahkan sebagai rumah makan nasi padang dengan rasa otentik, \textit{brand recognition} yang kuat, produk unggulan jelas, dan dukungan promo platform                        & Pemilik menekankan rasa otentik dan \textit{brand recognition}; observasi platform menunjukkan menu unggulan dan promo aktif \\ \hline
        \bottomrule
    
\end{xltabular}
\tablesource[\textwidth]{Wawancara awal dengan pemilik usaha; diolah peneliti (2026).}
\endgroup

\begingroup
\begin{xltabular}{\textwidth}{|V|C{3.0cm}|C{3.0cm}|C{2.0cm}|}
\caption{Komposisi Pelanggan Berdasarkan Keterangan Pemilik}
\label{tab:komposisi_pelanggan_wawancara} \\
        \toprule
        \textbf{Kategori Pelanggan} & \textbf{\shortstack{Periode Sebelumnya\\sebelum \textit{high season}}} & \textbf{\shortstack{Saat Wawancara Awal\\April 2026}} & \textbf{Perubahan} \\ \hline
        \midrule
\endfirsthead
\continuedtabletitle{4}
        \toprule
        \textbf{Kategori Pelanggan} & \textbf{\shortstack{Periode Sebelumnya\\sebelum \textit{high season}}} & \textbf{\shortstack{Saat Wawancara Awal\\April 2026}} & \textbf{Perubahan} \\ \hline
        \midrule
\endhead
\continuedtablefooter{4}
\endfoot
        \textit{New-comers}         & 45\%                        & 55\%              & +10 pp             \\ \hline
        \textit{Repeat order}       & 43\%                        & 38\%              & -5 pp              \\ \hline
        \textit{Reactivated user}   & 12\%                        & 7\%               & -5 pp              \\ \hline
        \bottomrule
    
\end{xltabular}
\tablesource[\textwidth]{Wawancara awal dengan pemilik usaha; diolah peneliti (2026).}
\endgroup
Tabel~\ref{tab:profil_wawancara_awal}, Tabel~\ref{tab:wawancara_platform}, Tabel~\ref{tab:stp_wawancara_awal}, dan Tabel~\ref{tab:komposisi_pelanggan_wawancara} menunjukkan bahwa wawancara awal memberikan data mengenai identitas usaha, alasan pemilihan platform utama, indikasi STP awal, serta perubahan komposisi pelanggan dari periode sebelum \textit{marketing plan high season} ke kondisi saat wawancara awal pada April 2026. Pemilik usaha mengaitkan perubahan komposisi pelanggan tersebut dengan pelaksanaan \textit{marketing plan high season}.

\begingroup
\begin{xltabular}{\textwidth}{|Z{3.4cm}|Y|}
\caption{Keterangan Pemilik Mengenai Kekuatan dan Strategi Usaha}
\label{tab:kekuatan_strategi_wawancara} \\
        \toprule
        \textbf{Aspek}            & \textbf{Keterangan Berdasarkan Wawancara}                                                                                                      \\ \hline
        \midrule
\endfirsthead
\continuedtabletitle{2}
        \toprule
        \textbf{Aspek}            & \textbf{Keterangan Berdasarkan Wawancara}                                                                                                      \\ \hline
        \midrule
\endhead
\continuedtablefooter{2}
\endfoot
        Kekuatan utama            & Rasa produk dinilai otentik dan \textit{brand recognition} dinilai sudah cukup kuat                                                            \\ \hline
        Produk unggulan           & Dendeng dan ayam pop putih                                                                                                                     \\ \hline
        Strategi                  & \textit{Upselling} melalui optimisasi \textit{product lineup}; peningkatan kualitas produk; peningkatan layanan; \textit{discount partnership} \\ \hline
        Peran GrabFood            & Digunakan lebih banyak sebagai kanal \textit{call to action}                                                                                   \\ \hline
        Kanal pemasaran utama     & Media sosial serta KOL atau \textit{influencer}                                                                                                \\ \hline
        Penguatan identitas merek & Perancangan maskot, \textit{signature gesture}, dan budaya merek yang unik                                                                     \\ \hline
        \bottomrule
    
\end{xltabular}
\tablesource[\textwidth]{Wawancara awal dengan pemilik usaha; diolah peneliti (2026).}
\endgroup

\begingroup
\begin{xltabular}{\textwidth}{|Z{3.4cm}|Y|}
\caption{Kendala Operasional Berdasarkan Wawancara Pemilik}
\label{tab:kendala_wawancara} \\
        \toprule
        \textbf{Kategori Kendala}       & \textbf{Keterangan Berdasarkan Wawancara}                                                                                                                                                                                                                         \\ \hline
        \midrule
\endfirsthead
\continuedtabletitle{2}
        \toprule
        \textbf{Kategori Kendala}       & \textbf{Keterangan Berdasarkan Wawancara}                                                                                                                                                                                                                         \\ \hline
        \midrule
\endhead
\continuedtablefooter{2}
\endfoot
        Kendala umum                    & Ongkos kirim dinilai mahal; jarak layanan masih terbatas; harga mengalami kenaikan karena usaha sedang meningkatkan margin dan \textit{net profit}; diskon, pajak, dan biaya promo tidak lagi sepenuhnya ditanggung platform                                      \\ \hline
        Kendala layanan                 & Antrean driver sering terjadi terutama pada \textit{peak season}; pemilik menginginkan sistem antrean yang memprioritaskan driver yang datang lebih dahulu; makanan yang habis masih dapat dipesan; pesanan kadang tidak sesuai dengan \textit{request} pelanggan \\ \hline
        \textit{Bottleneck} operasional & Pada \textit{peak hour} banyak menu \textit{sold out} dan perlu \textit{restock}; pelanggan dapat membatalkan pembelian ketika menu habis; kapasitas penggorengan masih kurang; belum tersedia opsi \textit{dessert} untuk mendukung \textit{upselling}           \\ \hline
        Konsistensi rasa                & Pemilik menegaskan konsistensi rasa masih menjadi masalah utama; sebagian \textit{cook} masih mengambil \textit{shortcut}; diperlukan pengawasan yang lebih ketat agar rasa tetap konsisten                                                                       \\ \hline
        \bottomrule
    
\end{xltabular}
\tablesource[\textwidth]{Wawancara awal dengan pemilik usaha; diolah peneliti (2026).}
\endgroup

\begingroup
\begin{xltabular}{\textwidth}{|Z{3.4cm}|Y|}
\caption{Peluang dan Ancaman Berdasarkan Wawancara Pemilik}
\label{tab:peluang_ancaman_wawancara} \\
        \toprule
        \textbf{Kategori}  & \textbf{Keterangan Berdasarkan Wawancara}                                                                                                                                                                                           \\ \hline
        \midrule
\endfirsthead
\continuedtabletitle{2}
        \toprule
        \textbf{Kategori}  & \textbf{Keterangan Berdasarkan Wawancara}                                                                                                                                                                                           \\ \hline
        \midrule
\endhead
\continuedtablefooter{2}
\endfoot
        Peluang platform   & Dukungan program GrabFood seperti awards, \textit{verified store}, dan \textit{voucher discount}; peluang \textit{upselling} melalui pengembangan produk; potensi peningkatan nilai transaksi melalui \textit{discount partnership} \\ \hline
        Peluang ekspansi   & Usaha belum memiliki cabang; pemilik ingin membuka cabang kedua pada tahun ini; terdapat peluang membuka cabang baru di Kota Medan; target jangka panjang adalah ekspansi ke Jakarta setelah memiliki lima cabang di Medan          \\ \hline
        Ancaman persaingan & Munculnya kompetitor baru dengan \textit{brand recognition} kuat di Kota Medan, seperti Pagi Sore dan Payakumbuah                                                                                                                   \\ \hline
        Ancaman platform   & Tingginya ongkos kirim; pergeseran beban promo kepada konsumen; sensitivitas pelanggan terhadap harga                                                                                                                               \\ \hline
        \bottomrule
    
\end{xltabular}
\tablesource[\textwidth]{Wawancara awal dengan pemilik usaha; diolah peneliti (2026).}
\endgroup

\begingroup
\begin{xltabular}{\textwidth}{|Z{3.4cm}|Y|}
\caption{Kategorisasi SWOT Awal Berdasarkan Wawancara Pemilik}
\label{tab:swot_wawancara_awal} \\
        \toprule
        \textbf{Kategori}      & \textbf{Data Awal}                                                                                                                                                                                                                                                                                           \\ \hline
        \midrule
\endfirsthead
\continuedtabletitle{2}
        \toprule
        \textbf{Kategori}      & \textbf{Data Awal}                                                                                                                                                                                                                                                                                           \\ \hline
        \midrule
\endhead
\continuedtablefooter{2}
\endfoot
        \textit{Strengths}     & Rasa produk dinilai otentik; \textit{brand recognition} dinilai cukup kuat; GrabFood disebut sebagai penyumbang lebih dari 50\% pendapatan penjualan; terdapat produk unggulan berupa dendeng dan ayam pop putih; usaha mendapat dukungan program platform                                                       \\ \hline
        \textit{Weaknesses}    & Proporsi \textit{repeat order} menurun; konsistensi rasa belum stabil; makanan yang habis masih dapat dipesan; pesanan kadang tidak sesuai; terjadi \textit{bottleneck} operasional saat \textit{peak hour}; kapasitas penggorengan masih terbatas; belum tersedia \textit{dessert} untuk \textit{upselling} \\ \hline
        \textit{Opportunities} & Terdapat peluang \textit{upselling}; peluang \textit{discount partnership}; dukungan fitur dan program promosi GrabFood; rencana penguatan identitas merek; peluang ekspansi cabang di Medan dan Jakarta                                                                                                     \\ \hline
        \textit{Threats}       & Munculnya kompetitor baru dengan \textit{brand recognition} kuat; ongkos kirim dinilai mahal; jarak pengantaran terbatas; kenaikan harga berpotensi memengaruhi pembelian; beban promo semakin banyak ditanggung pengguna                                                                                    \\ \hline
        \bottomrule
    
\end{xltabular}
\tablesource[\textwidth]{Diolah peneliti berdasarkan wawancara awal dengan pemilik usaha (2026).}
\endgroup
Berdasarkan wawancara awal, pemilik usaha menyampaikan data mengenai posisi GrabFood sebagai platform utama, komposisi pelanggan, kekuatan produk, strategi pemasaran, kendala operasional, peluang pengembangan, dan ancaman eksternal. Seluruh data tersebut digunakan sebagai bahan awal untuk mengidentifikasi faktor internal dan eksternal usaha pada analisis SWOT penelitian.

\subsection*{Konteks Pra-Survei dalam Analisis Bab IV}
Selain data penelitian utama, analisis juga menggunakan data pra-survei sebagai konteks awal. Data pra-survei tidak menggantikan data lapangan Juni--Juli 2026, tetapi membantu membaca pola awal tentang penjualan, perilaku pelanggan, dan posisi Nasi Gerilya pada kanal digital. Kelompok pra-survei disajikan menurut sumber data agar konteks awal penelitian dapat dibaca secara lebih runtut.

\begingroup
    \scriptsize
    \begin{xltabular}{\textwidth}{|C{0.7cm}|Y|L{4.4cm}|Z{1.8cm}|}
        \caption{Konteks Pra-Survei dalam Analisis Bab IV}
        \label{tab:konteks-prasurvei-bab4}\\
        \toprule
        \textbf{No.} & \textbf{Temuan Penting} & \textbf{Fungsi dan Peran Analisis} & \textbf{Kode Sumber} \\ \hline
        \midrule
        \endfirsthead
        \continuedtabletitle{4}
        \toprule
        \textbf{No.} & \textbf{Temuan Penting} & \textbf{Fungsi dan Peran Analisis} & \textbf{Kode Sumber} \\ \hline
        \midrule
        \endhead
        \continuedtablefooter{4}
        \endfoot
        \bottomrule
        \endlastfoot
        1 & GrabFood disebut sebagai kanal utama dengan estimasi kontribusi lebih dari 50\% penjualan; pelanggan baru naik, tetapi repeat order dan reactivated user turun; dendeng dan ayam pop putih muncul sebagai produk unggulan; promo berpengaruh tetapi harus dihitung dengan pajak, HPP, dan margin; masalah stok, antrean driver, konsistensi rasa, dan pesanan tidak sesuai request sudah muncul sejak pra-survei. & Menguatkan fokus penelitian pada GrabFood, isu loyalitas, kebutuhan hero menu, promo terukur, update stok, dan SOP peak hour; menjadi konteks awal untuk membaca kontribusi GrabFood, pertimbangan margin, dan pola masalah layanan. & K08.1 \\ \hline
        2 & Observasi 4 April 2026 menunjukkan rating 4,7 dari 4.361 penilaian, variasi menu, promo, add-on, lauk tambahan, sambal kemasan, kuah/minuman, dan menu premium. & Menjadi pembanding historis terhadap observasi 13 Juni 2026, terutama untuk reputasi digital, rentang produk, promo, dan peluang upselling; angka terbaru yang dipakai untuk temuan utama tetap observasi 13 Juni 2026. & K08.2 \\ \hline
        3 & Data kinerja awal menunjukkan penjualan bulanan mencapai puncak pada Maret 2025 lalu cenderung melemah sampai November 2025; indeks Nasi Gerilya turun dari 100,0 menjadi 70,4 pada Juli--Oktober 2025 atau berubah -29,6\%; AOV dan kunjungan naik pada sebagian periode tetapi penjualan dan indeks relatif melemah. & Menjelaskan urgensi strategi peningkatan penjualan, kebutuhan menjaga konversi, repeat order, harga, promo, dan pengalaman layanan; menjadi indikator historis untuk membaca urgensi strategi, bukan ukuran dampak setelah rekomendasi diterapkan. & K08.3 \\ \hline
        4 & Foto produk dan tampilan toko menunjukkan variasi menu, termasuk ayam pop putih dan nasi dendeng sebagai menu unggulan. & Mendukung analisis \textit{product} dan \textit{physical evidence}, terutama penguatan foto, deskripsi, menu hero, dan tampilan digital; foto produk membuktikan ketersediaan aset tampilan produk, bukan membuktikan menu paling laku. & K08.4 \\ \hline
    \end{xltabular}
\tablesource[\textwidth]{Diolah peneliti berdasarkan koding pra-survei (2026).}
\endgroup

Dengan demikian, data pra-survei dipakai sebagai konteks historis dan penguat pola, sedangkan keputusan analisis utama disandarkan pada data lapangan Juni--Juli 2026. Setelah wawancara pemilik 6 Juli 2026 diperoleh, beberapa tema yang semula hanya menjadi konteks awal dapat dibaca sebagai temuan yang diperkuat oleh informan kunci, terutama komposisi kanal penjualan, konsistensi rasa, pengendalian bahan, struktur leader, promo, dan pengembangan sistem. Setiap kelompok data pra-survei tetap dibaca sesuai sumbernya terlebih dahulu, kemudian hanya digunakan sebagai pembanding ketika temanya muncul kembali pada hasil wawancara, angket pelanggan, observasi, atau dokumentasi utama.
